\documentclass[11pt]{article}
\usepackage[margin=1in]{geometry}
\usepackage[T1]{fontenc}
\usepackage[utf8]{inputenc}
\usepackage{amsmath,amssymb,amsthm}
\usepackage[hidelinks]{hyperref}
\newtheorem{theorem}{Theorem}
\newtheorem{lemma}[theorem]{Lemma}
\newtheorem{corollary}[theorem]{Corollary}
\newtheorem{proposition}[theorem]{Proposition}
\newcommand{\Hom}{\operatorname{Hom}}
\newcommand{\Ann}{\operatorname{Ann}}
\newcommand{\Tr}{\operatorname{Tr}}
\newcommand{\Frac}{\operatorname{Frac}}
\newcommand{\Fitt}{\operatorname{Fitt}}
\newcommand{\Jac}{\operatorname{Jac}}
\newcommand{\wtensor}{\mathbin{\widehat\otimes}}
\title{Noether differents as formal Frobenius multipliers\\and completed Cartier product maps}
\author{Research working paper; independent review pending}
\date{9 October 2026}
\begin{document}
\maketitle
\begin{abstract}
We replace the single trace discriminant of draft391 by every element of
a Noether different. A separability idempotent and integral traces give a
multiplier independent of the Frobenius exponent, without a Cohen--Macaulay
assumption. Finite free root modules then give compatible maps on the actual
completed diagonal. We also prove an exact completed-tensor argument for
finite-module duals, which generates every Cartier map of a formal product
by pairs of Cartier maps. Under an explicitly stated minimal-test-ideal
framework this gives the product inclusion needed by a subadditivity route.
The identification with a primary formal Jacobian ideal retains its named
published source facts. Actual completed-diagonal restriction and the full
envelope theorem remain to be proved. Original P12.3 remains open.
\end{abstract}

\section{Notation and the exact scope}
Let $(R,\mathfrak m,k)$ be a complete Noetherian normal local domain of
characteristic $p>0$ with perfect residue field. Fix a coefficient field
$k\subset R$ inducing this residue field. Set $q=p^e$, $T=R^{1/q}$, using
the root ring in $\Frac(R)^{1/q}$ and its natural $R$-module structure.
All source multiplications below are in this root convention. Output
multiplication by $d\in R$ corresponds to source multiplication by $d^q$
in $F_*^eR$ notation.

A regular normalization here means an actual regular Noetherian subring
$A\subset R$ with $R$ finite over $A$. Put $K=\Frac(A)$ and
$L=\Frac(R)$. The principal lifting theorem assumes $L/K$ separable.
There is no assertion that a single arbitrary normalization has this
property or that its discriminant has primary support.

Write $B_A=R\otimes_A R$, let $\mu_A:B_A\to R$ be multiplication, and set
\[
 N_A=\mu_A\bigl(\Ann_{B_A}\ker\mu_A\bigr),\qquad
 K_A=\Fitt_0^R\Omega_{R/A}.
\]
The completed $k$-diagonal is
\[
 S=R\wtensor_k R,\quad \mu:S\to R,\quad I=\ker\mu.
\]
For each $e$, let $\mathcal D_e$ be the image in $\Hom_R(T,R)$ of the
$S$-linear maps $S^{1/q}\to S$ carrying the full kernel of root
multiplication to $I$. This is the full compatible-root-diagonal condition,
and not compatibility only with the ordinary extended ideal $IS^{1/q}$.
No assertion about the BIG test ideal of $\mathcal D$ is built into this
definition.

\section{Regular finite subrings have the needed root modules}
\begin{lemma}\label{regularbase}
Every regular normalization $A\subset R$ as above is complete local of
characteristic $p$, with perfect residue field. For every $q$,
$A^{1/q}$ is finite free as an $A$-module.
\end{lemma}
\begin{proof}
Finiteness makes $R$ integral over $A$. Lying over shows that $A$ has only
one maximal ideal $\mathfrak n$, the contraction of $\mathfrak m$.
The extension of residue fields $F=A/\mathfrak n\subset k$ is finite.
Also $\sqrt{\mathfrak nR}=\mathfrak m$, so the two adic topologies on $R$
are equivalent. Thus the $\mathfrak n$-adic completion of its finite
$A$-module is $R$ itself. Exactness of completion for finite modules over a
Noetherian ring applied to $A\subset R$ embeds $\widehat A$ into $R$.
It follows that $\widehat A$ is finite as an $A$-module. Its reduction
modulo $\mathfrak n$ is generated by the class of $1$. Nakayama therefore
makes $A\to\widehat A$ surjective, hence an isomorphism.

The field $F$ is perfect. Indeed all successive purely inseparable root
extensions of $F$ embed in the perfect finite extension $k$. If $F$ were
not perfect, adjoining successive roots of a non-$p$th power would give
extensions of unbounded degrees, contrary to $[k:F]<\infty$.
The equicharacteristic Cohen structure theorem for the complete regular
ring $A$ now gives $A\simeq F[[x_1,\ldots,x_d]]$, where $d=\dim R$.
Grouping formal exponents modulo $q$ gives
\[
 A^{1/q}=\bigoplus_{0\le\alpha_i<q} A x^{\alpha/q}.
\]
Coefficient roots exist in $F$, and the disjoint exponent classes prove
uniqueness. This proves finite freeness. The chosen coefficient field of
$A$ need not equal the chosen coefficient field of $R$.
\end{proof}
The backgrounds in this lemma are the exact finite-module completion
theorem, lying over, Nakayama, and Cohen structure for a complete regular
equicharacteristic local ring. No flatness of $R$ over $A$ is used.

\section{The entire Noether different gives root containment}
\begin{theorem}\label{rootcontainment}
Assume $L/K$ is separable. For every $d\in N_A$ and every $q$,
\[
 dR^{1/q}\subset A^{1/q}[R]\subset L^{1/q}.
\]
Multiplication identifies $P_q=A^{1/q}\otimes_A R$ with $A^{1/q}[R]$.
The module $P_q$ is finite free over $R$, even if $R$ is not
Cohen--Macaulay.
\end{theorem}
\begin{proof}
Choose $z=\sum_{i=1}^r r_i\otimes s_i$ in
$\Ann_{B_A}\ker\mu_A$ with $\mu_A(z)=d$. After localizing $A$ to $K$,
the finite domain $R\otimes_A K$ is the field $L$. The image of $z$ in
$L\otimes_K L$ annihilates the kernel of multiplication.

Because $L/K$ is finite separable, this annihilator is the one-dimensional
$L$-summand generated by its diagonal separability idempotent. The image
of $z$ is $d$ times that idempotent. Here is the exact trace consequence.
In a splitting field, for the distinguished embedding and every other
$K$-embedding $\sigma$ of $L$, one has respectively
\[
 \sum_i r_i s_i=d,\qquad \sum_i r_i\sigma(s_i)=0.
\]
Consequently, for $w\in L$,
\[
 \sum_i r_i\Tr_{L/K}(s_iw)=dw.
\]
This follows by expanding the trace as the sum of all separable embeddings;
it does not require $\Tr(1)\ne0$.

Separable $L/K$ is linearly disjoint from $K^{1/q}/K$.
Moreover $LK^{1/q}=L^{1/q}$: both fields have degree $[L:K]$ over
$K^{1/q}$ and the first embeds in the second. The trace identity therefore
base changes to
\[
 dw=\sum_i r_i\Tr_{L^{1/q}/K^{1/q}}(s_iw),\qquad w\in L^{1/q}.
\]
If $w\in R^{1/q}$, each $s_iw$ is integral over $A^{1/q}$. Its trace is
integral and belongs to $K^{1/q}$, hence lies in the normal ring $A^{1/q}$.
This proves the asserted containment with the same $d$ for every $q$.

By Lemma\ref{regularbase}, $P_q$ is finite free over $R$. As an
$A$-module it is a finite direct sum of the torsion-free module $R$.
Its generic multiplication map is the preceding field isomorphism.
Thus its kernel is zero, and its image is precisely $A^{1/q}[R]$.
\end{proof}
This theorem uses an element of $R$, not just a discriminant in $A$.
Summing Noether differents from different actual normalizations is allowed
once the following compatibility theorem is proved.

\section{Compatible maps on the actual completed diagonal}
\begin{theorem}\label{lift}
Under the hypotheses of Theorem\ref{rootcontainment},
\[
 N_A\Hom_R(R^{1/q},R)\subset\mathcal D_e
 \quad\text{for every }e.
\]
\end{theorem}
\begin{proof}
Fix $d\in N_A$ and $\varphi:T\to R$ $R$-linear. Multiplication by $d$
gives an $R$-linear map $T\to P_q$. Put $B=R\otimes_k T$.
The module $P_q\otimes_k T$ is finite free over $B$.
Multiplication $B\to T$ is a surjective $B$-linear map, with $T$ acted on
through multiplication. Hence the multiplication map
$P_q\otimes_k T\to T$ lifts to a $B$-linear map
$\Psi:P_q\otimes_k T\to B$.
Compose
\[
 T\otimes_k T\xrightarrow{d\otimes1}P_q\otimes_k T
 \xrightarrow{\Psi}R\otimes_k T
 \xrightarrow{1\otimes\varphi}R\otimes_k R.
\]
This is $R\otimes_k R$-linear. Multiplying in the target gives
$\varphi(dab)=d\varphi(ab)$ on a tensor $a\otimes b$.
It therefore descends on the full root diagonal to $d\varphi$.

All these maps preserve every power of
$\mathfrak m\otimes1+1\otimes\mathfrak m$, by module linearity.
They extend to the corresponding completions. A presentation
$R=k[[u_1,\ldots,u_n]]/J$ identifies
\[
 S=k[[u,v]]/(J(u),J(v)),\qquad
 S^{1/q}=T\wtensor_k T.
\]
For the second equality take roots in the formal presentation; perfectness
of $k$ identifies its two root coefficient fields with $k$.
The root-adic and original adic topologies are cofinal. This is an explicit
presentation identity, not an assumption that completion is exact over the
possibly non-Noetherian algebraic tensor ring.
The commuting multiplication identity extends by continuity and density.
An element of the kernel of root multiplication is consequently sent into
$I$. The completed map is the required member of $\mathcal D_e$.
\end{proof}

\begin{lemma}\label{kahler}
For every module-finite extension $A\subset R$ of commutative Noetherian rings,
$K_A\subset N_A$.
\end{lemma}
\begin{proof}
Use a finite polynomial presentation
$R=A[X_1,\ldots,X_r]/(f_1,\ldots,f_s)$, available because $A$ is
Noetherian and $R$ is finite. In $B_A$ the kernel of multiplication is
generated by $t_i=x_i\otimes1-1\otimes x_i$.
Write each difference of the two evaluations of $f_j$ as
$\sum_i a_{ji}t_i=0$ in $B_A$, with
$\mu_A(a_{ji})=\partial f_j/\partial X_i$ in $R$.
For any choice of $r$ relations, the resulting square matrix $C$ satisfies
$Ct=0$. Its adjugate gives $(\det C)t_i=0$ for every $i$.
Thus $\det C\in\Ann_{B_A}\ker\mu_A$, while its multiplication image is
the corresponding Jacobian minor. These minors generate
$\Fitt_0^R\Omega_{R/A}$. The case $r=0$ is the identity extension.
\end{proof}

\section{The primary-ideal corollary and its source boundary}
Define the normalization Jacobian
\[
 J_{\mathrm{norm}}=\sum_{A=k[[x_1,\ldots,x_d]]\subset R\text{ finite}} K_A,
\]
where the fixed coefficient field $k$ is used and $(x_1,\ldots,x_d)$
runs through systems of parameters. These are actual regular finite
subrings. The lifting arguments above also apply to other regular finite
subrings, but the primary corollary needs only this specified sum.
If $K_A\ne0$, localization gives $\Omega_{L/K}=0$, so $L/K$ is
separable. Indeed a finite field extension has zero relative differentials
exactly when it is separable. Terms with $K_A=0$ contribute nothing.
The preceding theorems and $R$-module linearity give the unconditional
algebraic conclusion
\[
 J_{\mathrm{norm}}\Hom_R(T,R)\subset\mathcal D_e
 \quad\text{for every }e.
\]
Only finitely many normalizations are needed to generate this sum, since
$R$ is Noetherian. There is no choice of one normalization working for all
generators.

For the next corollary the exact additional published source facts are:
\begin{enumerate}
\item For an equicharacteristic complete local ring
$R=k[[u_1,\ldots,u_n]]/(f_1,\ldots,f_s)$ of dimension $d$,
$J_{\mathrm{norm}}$ equals the ideal of $(n-d)$-minors of the formal
Jacobian matrix. Equivalently it is the $d$th Fitting ideal of continuous
differentials. This is the equality stated in Iyengar--Takahashi,
\emph{The Jacobian ideal of a commutative ring and annihilators of
cohomology}, Example3.3, referring to Wang's Lemma4.3.
The latter source uses precisely coefficient-field-based normalizations
in Definition4.1. Its complete statement and immediate proof have now
also been read on the official scanned pages597--598.
\item For an equidimensional complete equicharacteristic local ring with
perfect residue field, this formal Jacobian ideal defines the singular
locus. Liu, \emph{Annihilators and dimensions of the singularity category},
Corollary4.9 and its immediate proof state precisely this local case and
refer to its underlying formal Jacobian criterion.
\end{enumerate}
Their statements and immediate cited source context have been read from
retained author/published extracts and the actual Wang pixels just named.
Wang's proof uses prime avoidance to choose minimal generators such that
every $d$-subset is a system of parameters. This need not mean choosing
linear coefficients in an infinite residue field: generators can be
modified within $\mathfrak m^2$ while preserving their residue basis.
Indeed, for a finite set of incomparable nonmaximal primes, choose for
each prime an element of $\mathfrak m^2$ in all the other primes and
outside that prime. Adding the sum corresponding to the primes that
contain a given generator avoids all of them and preserves its residue
class. Applying this to the finitely many minimal primes of the previous
partial parameter ideals gives the required successive parameter choices.
Thus no infinite-field hypothesis is added here. The deeper formal
Jacobian-criterion proof tree is not declared wholly audited.
The corollary retains these
two exact conventional algebraic foundations; independent review must
check their full version and applicability. Normal domains are
equidimensional, and all their displayed hypotheses are satisfied here.

\begin{corollary}\label{primary}
Retain these two published identifications. If $R$ has an isolated
singularity, there exists $N$ such that
\[
 \mathfrak m^N\Hom_R(T,R)\subset\mathcal D_e
 \quad\text{for every }e.
\]
In the regular case $J_{\mathrm{norm}}=R$ and the same conclusion holds
with exponent zero.
\end{corollary}
\begin{proof}
In the singular isolated case the formal Jacobian has radical
$\mathfrak m$, hence contains a primary power. Apply the preceding
compatible-image inclusion. The exponent $N$ is independent of $e$.
\end{proof}
This addresses the compatible primary multiplier part of the formal
route under the stated published identifications. It says nothing by
itself about the test-ideal restriction or envelope continuity theorem.
No Cohen--Macaulay assumption from a singularity-category theorem has been
imported. In particular Dao--Takahashi's Theorem1.1 is not used outside
its expressly Cohen--Macaulay hypotheses.

\section{Exact completed tensor products of finite modules}
\begin{lemma}\label{exacttensor}
Let $k$ be a field and $R_1,R_2$ complete Noetherian local $k$-algebras
with residue field $k$. Completed tensor over $k$ is exact in each
variable on their finite modules with the adic topologies.
\end{lemma}
\begin{proof}
Such a module is the inverse limit of finite-dimensional $k$-vector
spaces. Its topology is Hausdorff and complete. By Artin--Rees, a short
exact sequence of finite modules has the induced submodule topology and
the quotient topology. Let $V^\circ$ denote continuous linear forms
$V\to k$, where $k$ has the discrete topology.
Continuous dual turns this short exact sequence into an exact sequence:
a form on a submodule factors through its intersection with some open
subspace, and extends by elementary linear algebra on the corresponding
finite-dimensional quotient. Kernel and quotient assertions are direct.
Evaluation identifies $V$ with the full algebraic dual
$(V^\circ)^*$, equipped with its finite-test topology, since both sides
are the same inverse limit of the finite-dimensional quotients.

The definition by pairs of finite quotients consequently gives
\[
 V\wtensor_k W=(V^\circ\otimes_k W^\circ)^*.
\]
Tensor and algebraic dual are exact over a field. Applying this identity
to the dual short exact sequence proves exactness of completed tensor.
No flatness of a coherent module over either local ring is assumed.
\end{proof}

\begin{proposition}\label{homfactor}
For finite $R$-modules $M,N$, with $S=R\wtensor_k R$,
\[
 \Hom_S(M\wtensor_k N,S)
 =\Hom_R(M,R)\wtensor_k\Hom_R(N,R)
\]
under the natural map. In particular every $S$-linear map
$S^{1/q}\to S$ is a finite $S$-linear combination of paired maps
$\varphi\wtensor\psi$ with $\varphi,\psi\in\Hom_R(T,R)$.
\end{proposition}
\begin{proof}
Choose finite presentations of $M$ and $N$.
The tensor product has generators consisting of pairs of their finite
generators, with two sets of relations supplied by the two presentation
matrices. A map to $S$ is an array in $S^{rs}$ satisfying both sets of
dual matrix equations. First imposing the $N$ equations identifies the
array space with
$R^r\wtensor_k\Hom_R(N,R)$ by Lemma\ref{exacttensor} applied to the
kernel of the finite dual presentation of $N$. Imposing the $M$ equations
then gives precisely
$\Hom_R(M,R)\wtensor_k\Hom_R(N,R)$ by the same lemma.
All these finite kernels carry their adic topologies, again by
Artin--Rees. Thus the identification is the asserted natural map.

The two dual modules are finite. Choose finite generating lists for them.
Their paired generators generate the completed tensor as a finite
$S$-module, by the completed finite presentations. Taking $M=N=T$
and the root-presentation identity proves the last assertion.
\end{proof}

The ring $S$ is Noetherian by its explicit formal presentation.
It is F-finite because $T$ is finite over $R$ and
$S^{1/q}=T\wtensor_k T$ is finite over $S$.
It is reduced: the inclusions $R\to R^{1/p}$ in both variables remain
injective under exact completed tensor, and their tensor is precisely
$S\to S^{1/p}$. This is injectivity of Frobenius. These conclusions do
not need normality of $S$ or a flatness assumption on $T$ over $R$.

\section{A conditional BIG product inclusion}
The following proposition explicitly assumes a minimal-test-ideal
framework. For a ring and a pair of ideals/exponents, suppose its BIG
test ideal exists and is the smallest ideal containing a nonzerodivisor
that is stable under every map
\[
 x\longmapsto\varphi\bigl((a x)^{1/q}\bigr),\qquad
 a\in\mathfrak a^{\lceil t(q-1)\rceil},
 \quad\varphi\in\Hom_R(R^{1/q},R),
\]
and analogously for a product of two ideal powers.
The ring and the pair must satisfy this existence/minimality assertion;
it is not an axiom granted for arbitrary Cartier subalgebras here.
\begin{proposition}\label{product}
Assume this framework for $(R,\mathfrak a^t)$,
$(R,\mathfrak b^s)$ and
$(S,(\mathfrak aS)^t(\mathfrak bS)^s)$, with the two ideals extended
from their respective factors. Then
\[
 \tau_S((\mathfrak aS)^t(\mathfrak bS)^s)
 \subset
 \tau_R(\mathfrak a^t)\wtensor_k\tau_R(\mathfrak b^s).
\]
\end{proposition}
\begin{proof}
Set $U=\tau_R(\mathfrak a^t)$, $V=\tau_R(\mathfrak b^s)$.
Exact completed tensor embeds $U\wtensor V$ as an ideal of $S$.
Choose nonzero $c\in U$, $d\in V$. Multiplication by $c$ and $d$ is
injective in the corresponding copies of the domain $R$. Exact tensor
proves multiplication by $c\otimes d$ is injective on $S$, so this ideal
contains a nonzerodivisor.

Fix $q$. Proposition\ref{homfactor} expresses every Cartier map of $S$
as a finite $S$-linear combination of paired maps. Its multipliers in the
two specified ideal powers are also generated by pairs of finite ideal
generators. On pure paired inputs, stability of $U$ and $V$ puts every
such image in $U\wtensor V$. The maps are continuous for the adic
topology: finite $S$-module linearity implies continuity on $S^{1/q}$,
whose root and original adic topologies are cofinal. Pure tensor sums
are dense and the target ideal is closed by Noetherian completeness.
Therefore the same inclusion holds on every completed input.
Minimality of the stated BIG test ideal proves the result.
\end{proof}
This argument does not prove the missing diagonal restriction theorem.
In particular one must still show that the image on $R$ of the actual
product test ideal is generically nonzero and stable under the relevant
restricted Cartier algebra, establish its BIG test ideal with the
correct convention, and compare it to the full test ideal using the
uniform compatible multipliers. No arbitrary-subalgebra existence or
finite-type-to-formal restriction theorem is silently used.

\section{What remains for the original problem}
This working paper supplies direct formal root containment, compatible
completed-diagonal maps, finite-module dual factorization, and the stated
conditional product inclusion. The primary corollary depends exactly on
the displayed published formal Jacobian facts, whose deeper reading
boundary is explicit. Its mathematical correctness is independently
unreviewed. Whole391 remains separately retained with its correction:
completion of normalized links and completed cones was already proved
in accepted G331; its envelope and Jensen transport are corollaries.

To finish the isolated formal route one must prove actual BIG diagonal
restriction/subadditivity with precise all-map conventions, then carry
out the entire model/global trace-generation/primary uniform estimate
on projective schemes over the complete local base. The compatible
primary ideal alone is not that uniform envelope theorem. General
nonisolated higher dimension, mixed or imperfect characteristic, other
excellent scopes, and general Jensen lower hemicontinuity remain open.
No new original solution, accepted paper or correctness grant is claimed.
The current catalogue remains16 solved original atoms,26 open,97 accepted
scoped papers. The protected P16.4 editor is unchanged.

\begin{thebibliography}{9}
\bibitem{IT} S. B. Iyengar and R. Takahashi,
\emph{The Jacobian ideal of a commutative ring and annihilators of
cohomology}, J. Algebra571 (2021),280--296,
DOI10.1016/j.jalgebra.2018.07.034. Complete author PDF
\url{https://www.math.nagoya-u.ac.jp/~takahashi/jac3.pdf}; definition,
Lemma3.1 and Example3.3 are the precise applied source units.
\bibitem{Liu} J. Liu, \emph{Annihilators and dimensions of the singularity
category}, Nagoya Math. J.250 (2023),533--548,
DOI10.1017/nmj.2022.45. Complete published author-hosted PDF
\url{https://jianliu.org/annihilator.pdf}; Corollary4.9 and its full
immediate proof supply the exact formal Jacobian criterion boundary.
\bibitem{Wang} H.-J. Wang, \emph{On the Fitting ideals in free resolutions},
Michigan Math. J.41 (1994),587--608,
DOI10.1307/mmj/1029005082. Official scanned PDF retained; OCR and pixel
reading coverage are recorded separately. Retention is not proof reading.
\bibitem{DT} H. Dao and R. Takahashi, \emph{Upper bounds for dimensions of
singularity categories}, C. R. Math. Acad. Sci. Paris353 (2015),297--301.
The expressly Cohen--Macaulay Theorem1.1 is a scope diagnostic, not a
premise for the non-Cohen--Macaulay lifting argument here.
\end{thebibliography}
\end{document}
